\section{Estudio de Losas por piso}

\subsection{Espesores de losa}

Se tomaron las 3 losas de mayor dimensión de cada piso, se obtuvo la realción $\epsilon$
, se interpoló el valor $\phi$ en base a la \textbf{tabla del apunte del curso} y se determinó el tipo de losa según tipo de apoyos. Con eso
se utilizó la siguiente formula para obtener el espesor de la losa:

\begin{equation}
    e \ge \frac{L_x \phi}{\lambda + r}
\end{equation}

Donde $L_x$ es la luz menor, $\phi$ es el coeficiente de esbeltez, $\lambda$ es 35 si es losa típica y 45 si es losa de techo, $r$ es el recubrimiento (2 cm) y $e$ es el espesor de la losa.

Según los espesores de las losas utilizadas, se definirá el espesor final como el mayor valor entre las 3 losas.

\subsubsection{Piso 1}

Se tomarón las losas 101, 110 y 111. Los espesores obtenidos fueron los siquientes:

\begin{table}[H]
    \centering
    \begin{tabular}{|c|c|c|c|c|c|c|}
        \hline
        Losa & Tipo Losa & $L_x$ (m)& $L_y$ (m) & $\epsilon$ & $\phi$ & $e$ (cm) \\
        \hline
        101 & 6 & 6.24 & 6.91 & 1.11 & 0.55 & 9.27 \\
        110 & 6 & 5.78 & 7.12 & 1.23 & 0.56 & 8.75 \\
        111 & 3a & 5.73 & 6.12 & 1.07 & 0.6 & 9.29 \\
        \hline
    \end{tabular}
    \caption{Espesores de losas del piso 1}
\end{table}

Por lo tanto, el espesor final del piso 1 es $e=10$ cm.

\subsubsection{Piso 2}

Se tomarón las losas 201, 202 y 208. Los espesores obtenidos fueron los siquientes:

\begin{table}[H]
    \centering
    \begin{tabular}{|c|c|c|c|c|c|c|}
        \hline
        Losa & Tipo Losa & $L_x$ (m)& $L_y$ (m) & $\epsilon$ & $\phi$ & $e$ (cm) \\
        \hline
        201 & 6 & 6.24 & 6.91 & 1.11 & 0.55 & 9.27 \\
        202 & 6 & 6.18 & 8.43 & 1.36 & 0.567 & 9.47 \\
        208 & 6 & 5.73 & 8.75 & 1.52 & 0.591 & 9.14 \\
        \hline
    \end{tabular}
    \caption{Espesores de losas del piso 2}
\end{table}

Por lo tanto, el espesor final del piso 2 es $e=10$ cm.

\subsubsection{Piso 3 al 21}

Se tomarón las losas 301, 302 y 308. Los espesores obtenidos fueron los siquientes:

\begin{table}[H]
    \centering
    \begin{tabular}{|c|c|c|c|c|c|c|}
        \hline
        Losa & Tipo Losa & $L_x$ (m)& $L_y$ (m) & $\epsilon$ & $\phi$ & $e$ (cm) \\
        \hline
        301 & 6 & 5.73 & 8.75 & 1.52 & 0.591 & 9.14 \\
        302 & 6 & 6.18 & 8.43 & 1.36 & 0.567 & 9.47 \\
        308 & 6 & 6.21 & 8.98 & 1.45 & 0.575 & 9.65 \\
        \hline
    \end{tabular}
    \caption{Espesores de losas del piso 3 al 21}
\end{table}

Por lo tanto, el espesor final del piso 3 al 21 es $e=10$ cm.

\subsubsection{Subterráneo}

Se tomarón las losas 008, 033 y 034. Los espesores obtenidos fueron los siquientes:

\begin{table}[H]
    \centering
    \begin{tabular}{|c|c|c|c|c|c|c|}
        \hline
        Losa & Tipo Losa & $L_x$ (m)& $L_y$ (m) & $\epsilon$ & $\phi$ & $e$ (cm) \\
        \hline
        008 & 6 & 6.24 & 6.91 & 1.11 & 0.55 & 9.27 \\
        033 & 6 & 5.78 & 7.12 & 1.23 & 0.56 & 8.75 \\
        034 & 2b & 5.60 & 8.15 & 1.45 & 0.945 & 14.30 \\
        \hline
    \end{tabular}
    \caption{Espesores de losas del subterráneo}
\end{table}

Por lo tanto, el espesor final del subterráneo es $e=15$ cm.
